
% 简洁风格：浅色背景 + 细边框 + 行号 + 自动换行
\definecolor{codebg}{RGB}{248,249,250}
\definecolor{codeframe}{RGB}{220,223,228}
\definecolor{codekw}{RGB}{25,85,145}
\definecolor{codecm}{RGB}{90,125,90}
\definecolor{codestr}{RGB}{136,44,136}

\lstdefinestyle{cleancpp}{
	language=C++,
	basicstyle=\ttfamily\small,
	keywordstyle=\bfseries\color{codekw},
    commentstyle=\color{codecm},
	stringstyle=\color{codestr},
	backgroundcolor=\color{codebg},
	frame=single,
	rulecolor=\color{codeframe},
	framesep=6pt,
	numbers=left,
	numberstyle=\scriptsize\color{gray},
	numbersep=8pt,
	breaklines=true,
	breakatwhitespace=false,
	showstringspaces=false,
	keepspaces=true,
	tabsize=4,
	columns=fullflexible
}

\chapter{Exercise 5.1}

\section{Problem}

Given the following grammar:

\[
S \to (L) \;|\; a
\]
\[
L \to L , S \;|\; S
\]

Tasks:
\begin{enumerate}
	\item Construct an LL(1) parsing table for the grammar. (Note: you must eliminate the left recursion first.)
	\item Draw the detailed process of the parsing of the sentence ``(a, (a, a))", follow the style in the previous slides.
\end{enumerate}

\section{Solution}

\subsection{Step 1: Eliminate left recursion}

The nonterminal $L$ is left-recursive. We rewrite $L$ as:

\[
S \to (L) \;|\; a
\]
\[
L \to S L'
\]
\[
L' \to , S L' \;|\; \epsilon
\]

\subsection{Step 2: FIRST and FOLLOW sets}

\begin{itemize}
    \item $\text{FIRST}(S) = \{ (, a \}$
    \item $\text{FIRST}(L) = \text{FIRST}(S) = \{ (, a \}$
    \item $\text{FIRST}(L') = \{ ,, \epsilon \}$
    \item $\text{FOLLOW}(S) = \{ ,, ), \$ \}$ (Note: $S$ is followed by $L'$ which starts with ',' or ends $L$, and $L$ is followed by ')')
    \item $\text{FOLLOW}(L) = \{ ) \}$
    \item $\text{FOLLOW}(L') = \text{FOLLOW}(L) = \{ ) \}$
\end{itemize}

\subsection{Step 3: LL(1) Parsing Table}

\begin{table}[htbp]
\centering
\begin{tabular}{|c|c|c|c|c|c|}
\hline
Nonterminal & ( & ) & a & , & \$ \\ \hline
$S$ & $S \to (L)$ & & $S \to a$ & & \\ \hline
$L$ & $L \to SL'$ & & $L \to SL'$ & & \\ \hline
$L'$ & & $L' \to \epsilon$ & & $L' \to ,SL'$ & \\ \hline
\end{tabular}
\end{table}


\subsection{Step 4: Parsing process for ``(a, (a, a))"}

\begin{table}[htbp]
\centering
\begin{tabular}{|l|r|l|l|}
\hline
Stack & Input & Action & Output \\\hline
\$S & (a, (a, a)) & derive & $S \to (L)$  \\
\$ ) L ( & (a, (a, a)) & match ( & \\
\$ ) L & a, (a, a)) & derive & $L \to SL'$ \\
\$ ) L' S & a, (a, a)) & derive & $S \to a$ \\
\$ ) L' a & a, (a, a)) & match a & \\
\$ ) L' & , (a, a)) & derive & $L' \to ,SL'$ \\
\$ ) L' S , & , (a, a)) & match , & \\
\$ ) L' S & (a, a)) & derive & $S \to (L)$  \\
\$ ) L' ) L ( & (a, a)) & match ( & \\
\$ ) L' ) L & a, a)) & derive & $L \to SL'$ \\
\$ ) L' ) L' S & a, a)) & derive & $S \to a$ \\
\$ ) L' ) L' a & a, a)) & match a & \\
\$ ) L' ) L' & , a)) & derive & $L' \to ,SL'$ \\
\$ ) L' ) L' S , & , a)) & match , & \\
\$ ) L' ) L' S & a)) & derive & $S \to a$ \\
\$ ) L' ) L' a & a)) & match a & \\
\$ ) L' ) L' & )) & derive & $L' \to \epsilon$ \\
\$ ) L' ) & )) & match ) & \\
\$ ) L' & ) & derive & $L' \to \epsilon$  \\
\$ ) & ) & match ) & \\\hline
\end{tabular}
\end{table}

\chapter{Exercise 5.2**}

\section{Problem}

Given the following grammar:

\[
A \to B \;|\; BC
\]
\[
B \to a B \;|\; \epsilon
\]
\[
C \to a b
\]

Tasks:
\begin{enumerate}
    \item Left factor the grammar.
    \item After left factoring, is the grammar an LL(1) grammar? or is it an LL(k) grammar? and why? (Note: you may try the input string ``ab".)
\end{enumerate}

\section{Solution}

\subsection{Step 1: Left factoring}

The production $A \to B \;|\; BC$ has a common prefix $B$. We factor it as:

\[
A \to B A'
\]
\[
A' \to C \;|\; \epsilon
\]
\[
B \to a B \;|\; \epsilon
\]
\[
C \to a b
\]

\subsection{Step 2: Analysis}

Grammar after left-factoring:
\[
A\to B A',\quad A'\to C\mid\epsilon,\quad B\to aB\mid\epsilon,\quad C\to ab.
\]

Proof:

\textbf{Not LL(1): }

Consider the input string ``ab''. On seeing the first symbol `a' we cannot decide with only one-symbol lookahead whether to use $B\to aB$ or $B\to\epsilon$: if $B\to\epsilon$ then $A\Rightarrow B A'\Rightarrow C\Rightarrow ab$ matches ``ab'', but if $B\to aB$ the sentential form will contain an extra `a' and cannot produce ``ab''. Thus a FIRST-set conflict exists, so the grammar is not LL(1).

\textbf{LL(2): }

With two-symbol lookahead the choice is unambiguous: if the next two symbols are ``ab'' choose $B\to\epsilon$ (so $A'\to C\to ab$); if they are ``aa'' choose $B\to aB$ to consume the extra `a'. Therefore a predictive parser with $k=2$ has no conflicts, so the grammar is LL(2).

Conclusion: the grammar is not LL(1) but is LL(2).

